% !TEX root = ../main.tex
\section{Illustrative measurement design}
\label{sec:methods}

\subsection{Endpoint and denominator}
The example asks how often an assigned trial reaches a prespecified target
verdict. A real experiment must define that target before execution and justify
its relationship to the research question. Here, it is simply the illustrative
endpoint used to explain the rate.

We need a denominator that retains unsuccessful assigned trials. For each
treatment $t\in\{A,B\}$, let $n_t>0$ be the number of assigned trials, let
$k_t\in\{0,\ldots,n_t\}$ count trials reaching the target, and let $f_t$ count
infrastructure failures. Failed trials do not enter $k_t$, but remain in $n_t$.
The assigned-trial target-verdict rate and descriptive contrast are
\begin{align}
  \assignedrate_t &= \frac{k_t}{n_t},
  \label{eq:assigned-rate}\\
  \ratecontrast &= \assignedrate_B - \assignedrate_A.
  \label{eq:contrast}
\end{align}
The rate is dimensionless; multiplying it by 100 expresses it as a percentage.
The corresponding contrast is reported in percentage points. When no trials
are assigned, the rate is undefined rather than zero.

The hypothetical setup contains 40 independently initialized trials per
treatment, with no shared memory. Its invented counts have no associated model
calls, executed backend, run identifiers, or deployment identity.

\subsection{Uncertainty and interpretation}
The figure uses marginal 95\% Wilson score intervals for each rate. These
intervals illustrate uncertainty under an independent binomial-trial model.
They do not establish that the hypothetical independence assumption holds in
a real agent experiment. Repeated episodes, shared scenarios, or persistent
agent memory may require a different analysis.

The intervals describe individual rates. Their overlap is not a test of the
treatment contrast in Equation~\eqref{eq:contrast}. A real comparison would
need an analysis of that contrast consistent with the assignment procedure,
sampling unit, and dependence structure.
